\documentclass[11pt]{article}

\usepackage[margin=1in]{geometry}
\usepackage[T1]{fontenc}
\usepackage[utf8]{inputenc}
\usepackage{setspace}
\usepackage[hyphens]{url}
\usepackage{hyperref}
\usepackage{xurl}

\setlength{\parindent}{0pt}
\setlength{\parskip}{6pt}

\begin{document}

\title{Group Assignment 5}
\author{Logan Rechin, Joe White\\ECO 442}
\date{\today}
\maketitle

\section*{The Plaza Accord and the Appreciation of the Yen:}

One event to research would be the Plaza Accord in September 1985, in which the G5 countries agreed to take action to weaken the US dollar. This event would be particularly interesting in the context of Japan because the impact on the exchange rate was so large. IMF material on the Plaza Accord notes that the Japanese yen ``appreciated by about 46 percent against the US dollar by the end of 1986.'' A paper by the Bank of Japan on the Plaza Accord and the exchange rate history shows the dollar fell to 211 on October 4, 1985, down from 242 just before the Plaza Accord. I think this event would be a good case to research because I find the connection to the exchange rate and the economic impact interesting. I also think it would be a good case to choose for the final project because the date is clear, the impact on the currency was large, and there was much discussion on the impact on the Japanese economy.

\textbf{Source:} \url{https://www.imf.org/-/media/websites/imf/imported-flagship-issues/external/pubs/ft/weo/2011/01/c1/_box14pdf.pdf}

\section*{Japan's Yen-Buying Intervention on September 22, 2022}

Another good event is the decision of the Japanese government to intervene in the foreign exchange markets on September 22, 2022. The Reuters news agency reports the following about this event: ``This was the first intervention in the foreign exchange markets by Japan in ten years. The intervention was carried out just after the Bank of Japan decided to maintain its ultra-low interest rate policy even as other countries were raising interest rates. The dollar fell more than 2 percent versus the yen after the intervention was announced.'' I think this event is interesting because it is a very pure example of the government trying to fight the rapid depreciation of its currency. This might actually be the easiest event for your final report because it occurred on one specific day and has a very clear exchange rate movement.

\textbf{Source:} \url{https://www.reuters.com/markets/asia/boj-keep-ultra-low-rates-remain-global-outlier-despite-weak-yen-2022-09-21/}

\section*{The Bank of Japan ends Negative Interest Rates}

The third interesting event is the decision by the Bank of Japan in March 2024 to remove negative interest rates. Reuters reported it as the end of eight years of negative interest rates and Japan's first rate hike in 17 years, as well as the abandonment of its yield curve control. What makes this event interesting is the fact that the Japanese yen weakened against the dollar after the event, even going below 150 per dollar, as the markets saw the guidance from the Bank of Japan as very dovish. This makes it a very interesting event from an economic perspective, as it shows that exchange rate effects are not just dependent on the movement in interest rates but also on expectations regarding interest rate differentials with the United States. This would be a great option for your final project if you and your group want something directly involving monetary policy.

\textbf{Source:} \url{https://www.reuters.com/markets/asia/japan-poised-end-negative-rates-closing-era-radical-policy-2024-03-18/}

\end{document}